\documentclass[amsmath, amssymb, aps]{revtex4-2}

\usepackage{graphicx}
\usepackage{dcolumn}
\usepackage{bm}
\usepackage{booktabs}
\usepackage{hyperref}
\usepackage{algorithm}
\usepackage{algpseudocode}
\usepackage{caption} 
\captionsetup[algorithm]{labelfont=bf, textfont=bf}
\usepackage{float}
\usepackage{placeins}
\usepackage{geometry}
\usepackage{setspace}
\usepackage{etoolbox}
\makeatletter
\@removefromreset{equation}{section}
\makeatother
\geometry{margin=1in}
\onehalfspacing

\hypersetup{
    colorlinks=true,
    linkcolor=blue,
    filecolor=magenta,      
    urlcolor=cyan,
}


\setlength{\parskip}{1em}

\begin{document}

\pagestyle{plain}

%\preprint{APS/123-QED}

\title{Four New Multi-Phase Hybrid Bracketing Algorithms for Numerical Root Finding }

\author{Abdelrahman Ellithy$^{1}$}
\affiliation{Department of Artificial Intelligence, Faculty of Computers and Artificial Intelligence, Benha University, Benha, Egypt}

\date{\today}


\begin{abstract}
Finding the roots of nonlinear equations requires balancing the efficiency of sophisticated techniques with the guaranteed convergence of traditional bracketing approaches. This manuscript presents four innovative hybrid algorithms: \textbf{Optimized Bisection-False Position (Opt.BF)}, \textbf{Optimized Bisection-False Position-Modified Secant (Opt.BFMS)}, \textbf{Optimized Trisection-False Position (Opt.TF)}, and \textbf{Optimized Trisection-False Position-Modified Secant (Opt.TFMS)}. These methods systematically integrate dependable interval contraction with accelerated interpolation processes. We analyze their convergence characteristics theoretically and evaluate them on a comprehensive benchmark of 48 diverse problems, including polynomial, transcendental, and ill-conditioned functions. A rigorous comparative analysis is performed against traditional benchmarks (Bisection, False Position) and established high-performance techniques (Brent's method, Quadratic Interval Refinement). The results indicate that the superior algorithm, \textbf{Opt.BFMS}, attains the most favorable overall performance, excelling in \textbf{66.7\%} of the problems and exhibiting substantial acceleration over both the quickest traditional baseline (Bisection, \textbf{6.0$\times$}) and leading-edge algorithms such as Brent's method (\textbf{1.75$\times$}) and the QIR method (\textbf{2.46$\times$}) on average. These findings substantiate the proposed hybrid methodologies as robust, efficient, and theoretically sound alternatives for numerical root finding.
\end{abstract}

\maketitle

\noindent\textbf{Keywords:} hybrid root-finding algorithms, root finding, Continuous optimization, convergence analysis, Derivative-free methods, Bracketing methods, modified secant method, Nonlinear equations, numerical analysis, numerical stability.

\section{Introduction}

Finding roots is a basic problem in numerical analysis that is used in many fields, including physics, engineering, economics, and computer science. Convergence is guaranteed with traditional root bracketing methods like Bisection and Regula Falsi (False Position), but they often have slow convergence rates, and on the other hand, the open methods are faster but sometimes diverge. This trade-off between reliability and speed has led to the creation of hybrid methods that preserve the strength of bracketing methods while converging faster with multi-phase accelerated interpolation techniques. This strategy has shown considerable promise in earlier work \cite{sabharwal2019blended, badr2022novel}. 

The field of root-finding algorithms has become more advanced beyond these traditional methods. Modern algorithms such as the widely used Brent's method \cite{brent1973algorithms} and Quadratic Interval Refinement (QIR) \cite{abbott2014quadratic} are now considered benchmarks for high performance. Even though these techniques are successful, there is still a chance to create new algorithms that are either quicker or perform better. This paper addresses that gap by introducing four new hybrid algorithms and evaluating them not only against traditional baselines but also against these established, high-performance algorithms.

The main contributions of this paper are:
\begin{enumerate}
    \item The design and implementation of \textbf{four novel hybrid algorithms} that blend reliable interval reduction utilize Bisection, Trisection and Regula Falsi with quasi-Newton steps.
    \item A \textbf{rigorous theoretical analysis} of the convergence order and error bounds for each proposed algorithm.
    \item A \textbf{comprehensive public presentation evaluation} on a diverse benchmark of 48 test functions, augmented with notoriously difficult cases to trial robustness and numerical stability.
    \item A \textbf{comparative analysis} against both traditional baselines and modern cutting-edge methods, using metrics that include function evaluations and performance in finite-precision arithmetic.
    \item The empirical finding that our premier hybrid algorithm, \textbf{Opt.BFMS}, provides a superior balance of speed and reliability, demonstrably outperforming the established benchmarks on average.
\end{enumerate}

This paper is organized as follows. Section~\ref{sec:related_work} (page~\pageref{sec:related_work}) reviews \textbf{related work} and contextualizes our contribution. Section~\ref{sec:methodology} (page~\pageref{sec:methodology}) details the \textbf{methodology} and mathematical foundations of the proposed algorithms. Section~\ref{sec:theoretical_analysis} (page~\pageref{sec:theoretical_analysis}) provides a \textbf{Theoretical Analysis} of their convergence properties. Section~\ref{sec:experimental} (page~\pageref{sec:experimental}) describes the \textbf{experimental setup} and test functions. Section~\ref{sec:results} (page~\pageref{sec:results}) presents the \textbf{results and analysis}. Section~\ref{sec:discussion} (page~\pageref{sec:discussion}) discusses the implications. Section~\ref{sec:conclusion} (page~\pageref{sec:conclusion}) concludes the paper and outlines future research directions. Section~\ref{sec:declaration} (page~\pageref{sec:declaration}) provides the \textbf{declarations}, including funding and competing interests statements.

\section{Related Work and Context}\label{sec:related_work}

The concept of hybridizing root-finding methods to achieve both reliability and speed is well-established, dating back to the influential algorithm of the 1970s. Dekker's algorithm ( 1969 ) was a seminal contribution, combining the safety of bisection with the velocity of the secant method. It introduced the now-common strategy of falling back on a bisection step if the faster insertion measure fails to make sufficient progress. Brent's method acting \cite{brent1973algorithms} refined this idea by replacing the secant method with inverse quadratic interpolation, achieving a higher order of convergence while retaining the bisection safeguard, making it one of the most widely used and powerful root-finders today.

Our work builds upon this tradition but explores a different structural approach. As a conditional fallback, our algorithm applies a deterministic, multi-phase sequence inside each loop, whereas Dekker's and Brents' method uses bisection. Think about it this way: for example, the Opt.BFMS algorithmic rule systematically applies a bisection, a false position, and a modified secant step, leveraging the output of each to feed the next as input. Avoiding the complex acceptance criteria of classic hybrids that aim to maximize progress in every iteration.

More recently, Quadratic Interval Refinement (QIR) \cite{abbott2014quadratic} has been proposed as a method that guarantees quadratic polynomial convergence for simple roots by fitting a quadratic model to three points to refine the root's location. While powerful, its performance can depend on the function's local behavior. The algorithms presented in this paper are differentiated by their reliance on simpler linear insertion and finite-difference steps, which, as our experiments show, can be more efficient on average across a wider variety of mapping types. 

Furthermore, this work is distinct from other recent hybrid methods included in our benchmark, such as those by Sabharwal (2019) \cite{sabharwal2019blended} and Badr (2021) \cite{badr2021comparative}. Those methods typically employ a blended strategy, often involving a conditional selection between outputs, for example, a bisection or a false position update within an iteration. Certainly, our advice algorithm differs fundamentally in its deterministic, \textbf{sequential multi-phase architecture}. For instance, our Opt.BF algorithm doesn't choose between bisection and false position. It executes a bisection, feeds the new output intervals as input to regula falsi. This fixed sequential process is a unique architectural choice designed to enforce procession while consistently applying interpolation-based acceleration.

\section{Algorithm Design and Methodology}\label{sec:methodology}

This study introduces four hybrid root-finding algorithms that share a unified, deterministic multi-phase architecture. Unlike traditional hybrid methods (e.g., Brent, Dekker) that employ complex conditional logic to switch between methods, our approach executes a fixed sequence of distinct phases within each iteration. This design systematically leverages the reliability of interval contraction and the convergence speed of interpolation without the overhead of "choose-the-best" decision trees.

\subsection{General Multi-Phase Architecture}
Each iteration $k$ of the proposed algorithms transforms the current bracketing interval $[a_k, b_k]$ into a refined interval $[a_{k+1}, b_{k+1}]$ through a cascade of up to three phases:
\begin{enumerate}
    \item \textbf{Phase 1: Guaranteed Reduction.} A bracketing step (Bisection or Trisection) ensures the search space contracts by a fixed factor ($1/2$ or $1/3$). This guarantees global convergence.
    \item \textbf{Phase 2: Linear Interpolation.} A False Position (Regula Falsi) step is applied to the interval produced by Phase 1. This accelerates convergence when the function behaves linearly.
    \item \textbf{Phase 3: Quadratic Acceleration (Optional).} A Modified Secant step estimates the root using a finite-difference approximation of the derivative. This step is accepted only if strictly robust safeguards are met.
\end{enumerate}

\subsection{Computational Building Blocks}
The algorithms are constructed from the following fundamental operators:

\textbf{Operator 1: Interval Halving (Bisection).}
Given $[a, b]$, computes $m = (a+b)/2$. The sub-interval containing the root is retained.

\textbf{Operator 2: Interval Trisection.}
Divides $[a, b]$ into three equal segments using $x_1 = a + (b-a)/3$ and $x_2 = b - (b-a)/3$. The algorithm evaluates $f(x_1)$ and $f(x_2)$ to identify the specific third containing the root, offering a theoretical reduction factor of $3$ per phase at the cost of two function evaluations.

\textbf{Operator 3: Regula Falsi (False Position).}
Computes the linear interpolation root $x_{FP}$:
\begin{equation}
    x_{FP} = \frac{a \cdot f(b) - b \cdot f(a)}{f(b) - f(a)}
\end{equation}
To prevent division by zero, a machine-epsilon guard ($\epsilon = 10^{-20}$) is applied to the denominator.

\textbf{Operator 4: Modified Secant.}
Approximates Newton's method using a finite difference step $\delta$. The candidate $x_S$ is computed from the current estimate $x_{FP}$:
\begin{equation}
    x_S = x_{FP} - \frac{f(x_{FP}) \cdot \delta}{f(x_{FP} + \delta) - f(x_{FP})}
\end{equation}

In this study, we employ an adaptive perturbation $$ \delta = 10^{-8} \times \max(1.0, |x_{FP}|) + 10^{-15} $$ to balance truncation and round-off errors, as recommended for robust numerical stability in derivative-free methods. This adaptive approach enhances generality across different root scales while maintaining performance in our benchmark. A candidate $$ x_S $$ is accepted if and only if it resides strictly within the current bracket and yields a smaller function magnitude than $$ x_{curr} $$.

\subsection{Algorithm Design, Hybrid Techniques and Pseudocode for Proposed Algorithms}

The four algorithms are distinguished by their specific combination of these phases:

\textbf{1. Optimized Bisection-False Position (Opt.BF).}
A two-phase sequence:
\begin{itemize}
    \item \textit{Step 1:} Execute Bisection.
    \item \textit{Step 2:} Execute Regula Falsi on the result of Step 1.
\end{itemize}

\textbf{2. Optimized Bisection-False Position-Modified Secant (Opt.BFMS).}
A three-phase sequence designed for maximum convergence speed:
\begin{itemize}
    \item \textit{Step 1:} Execute Bisection.
    \item \textit{Step 2:} Execute Regula Falsi.
    \item \textit{Step 3:} Execute Modified Secant on the result of Step 2, subject to acceptance criteria.
\end{itemize}

\textbf{3. Optimized Trisection-False Position (Opt.TF).}
Replaces binary splitting with ternary splitting:
\begin{itemize}
    \item \textit{Step 1:} Execute Trisection.
    \item \textit{Step 2:} Execute Regula Falsi on the chosen third.
\end{itemize}

\textbf{4. Optimized Trisection-False Position-Modified Secant (Opt.TFMS).}
The most aggressive variant:
\begin{itemize}
    \item \textit{Step 1:} Execute Trisection.
    \item \textit{Step 2:} Execute Regula Falsi.
    \item \textit{Step 3:} Execute Modified Secant on the result of Step 2.
\end{itemize}

\subsubsection{Algorithm 1: Optimized Bisection-False Position (Opt.BF)}

\captionof{algorithm}{Pseudocode for the Optimized Bisection-False Position (Opt.BF) algorithm. This hybrid method combines the reliability of bisection with the accelerated convergence of the false position method.}
    \begin{algorithmic}[1]
    \State \textbf{Input:} Function $f$, interval $[a, b]$, tolerance $tol$, max iterations $max\_iter$
    \State \textbf{Output:} Iteration count $n$, root approximation $x$, function value $f(x)$, final interval $[a, b]$
    \State $f_a \gets f(a)$, $f_b \gets f(b)$, $\epsilon \gets 10^{-20}$
    \State \textit{// Check if either bound is a root}
    \If{$|f_a| \leq tol$}
        \State \Return $1, a, f_a, a, b$
    \EndIf
    \If{$|f_b| \leq tol$}
        \State \Return $1, b, f_b, a, b$
    \EndIf
    \If{$f_a \times f_b \geq 0$}
        \State \Return $\text{None}, \text{None}, \text{None}, \text{None}, \text{None}$
    \EndIf
    \For{$n \gets 1$ \textbf{to} $max\_iter$}
        \State \textit{// Bisection Step}
        \State $mid \gets 0.5 \times (a + b)$
        \State $f_{mid} \gets f(mid)$
        \If{$|f_{mid}| \leq tol$}
            \State \Return $n, mid, f_{mid}, a, b$
        \EndIf
        \If{$f_a \times f_{mid} < 0$}
            \State $b, f_b \gets mid, f_{mid}$
        \Else
            \State $a, f_a \gets mid, f_{mid}$
        \EndIf
    
        \State \textit{// False Position Step with exception handling}
        \State $dx \gets (a \times f_b) - (b \times f_a)$
        \State \textbf{try:}
        \State \quad $fp \gets dx / (f_b - f_a)$
        \State \textbf{catch} (division by zero, overflow, value error):
        \State \quad $fp \gets dx / ((f_b - f_a) + \epsilon)$
        \State $f_{fp} \gets f(fp)$
        \If{$|f_{fp}| \leq tol$}
            \State \Return $n, fp, f_{fp}, a, b$
        \EndIf
        \If{$f_a \times f_{fp} < 0$}
            \State $b, f_b \gets fp, f_{fp}$
        \Else
            \State $a, f_a \gets fp, f_{fp}$
        \EndIf
    \EndFor
    \State \textit{// Max iterations reached}
    \State $final\_x \gets 0.5 \times (a + b)$
    \State $f_{final\_x} \gets f(final\_x)$
    \State \Return $max\_iter, final\_x, f_{final\_x}, a, b$
    \end{algorithmic}

\subsubsection{Algorithm 2: Optimized Bisection-False Position with Modified Secant (Opt.BFMS)}

The key step of this algorithm lies in its acceptance criterion for the modified secant estimate. The algorithm only accepts the secant output $x_S$ if it satisfies two conditions: (1) it lies within the current bracketing interval $[a, b]$, ensuring the bracketing property is maintained, and (2) it provides a better approximation than the false position estimate, i.e., $|f(x_S)| < |f(fp)|$. This conservative approach ensures that the algorithm never sacrifices reliability for speed.

\captionof{algorithm}{Pseudocode for the Hybrid Bisection-False Position-Modified Secant (HbisectionFalseMS) algorithm. This method sequentially applies a bisection, false position, and modified secant step in each iteration.}
    \begin{algorithmic}[1]
    \State \textbf{Input:} Function $f$, interval $[a, b]$, tolerance $tol$, max iterations $max\_iter$
    \State \textbf{Output:} Iteration count $n$, root approximation $x$, function value $f(x)$, final interval $[a, b]$
    \State $f_a \gets f(a)$, $f_b \gets f(b)$, $n \gets 0$, $\epsilon \gets 10^{-20}$
    \State \textit{// Check if either bound is a root}
    \If{$|f_a| \leq tol$}
        \State \Return $1, a, f_a, a, b$
    \EndIf
    \If{$|f_b| \leq tol$}
        \State \Return $1, b, f_b, a, b$
    \EndIf
    \If{$f_a \times f_b \geq 0$}
        \State \Return $\text{None}, \text{None}, \text{None}, \text{None}, \text{None}$
    \EndIf
    \While{$n < max\_iter$}
        \State $n \gets n + 1$
        
        \State \textit{// 1. Bisection Step to narrow the interval}
        \State $x_B \gets (a + b) \times 0.5$
        \State $f_{x_B} \gets f(x_B)$
        \If{$f_a \times f_{x_B} < 0$}
            \State $b \gets x_B$, $f_b \gets f_{x_B}$
        \Else
            \State $a \gets x_B$, $f_a \gets f_{x_B}$
        \EndIf
    
        \State \textit{// 2. False Position Step with exception handling}
        \State $dx \gets (a \times f_b) - (b \times f_a)$
        \State \textbf{try:}
        \State \quad $x_F \gets dx / (f_b - f_a)$
        \State \textbf{catch} (division by zero, overflow, value error):
        \State \quad $x_F \gets dx / ((f_b - f_a) + \epsilon)$
        \State $f_{x_F} \gets f(x_F)$
        \If{$f_a \times f_{x_F} < 0$}
            \State $b \gets x_F$, $f_b \gets f_{x_F}$
        \Else
            \State $a \gets x_F$, $f_a \gets f_{x_F}$
        \EndIf
        \If{$|f_{x_F}| \leq tol$}
            \State \Return $n, x_F, f_{x_F}, a, b$
        \EndIf
    
        \State \textit{// 3. Modified Secant Step with exception handling}
        \State $\delta \gets 10^{-8} \max(1, |x_F|) + 10^{-15}$
        \State \textbf{try:}
        \State \quad $x_S \gets x_F - \delta \times f_{x_F} / (f(x_F + \delta) - f_{x_F})$
        \State \textbf{catch} (division by zero, overflow, value error):
        \State \quad $x_S \gets x_F - \delta \times f_{x_F} / ((f(x_F + \delta) - f_{x_F}) + \epsilon)$
        \If{$a < x_S < b$}
            \State $f_{x_S} \gets f(x_S)$
            \If{$|f_{x_S}| < |f_{x_F}|$}
                \If{$f_a \times f_{x_S} < 0$}
                    \State $b \gets x_S$, $f_b \gets f_{x_S}$
                \Else
                    \State $a \gets x_S$, $f_a \gets f_{x_S}$
                \EndIf
                \If{$|f_{x_S}| \leq tol$}
                    \State \Return $n, x_S, f_{x_S}, a, b$
                \EndIf
            \EndIf
        \EndIf
    \EndWhile
    \State \textit{// Max iterations reached, return the midpoint of the final interval}
    \State $x \gets (a + b) \times 0.5$
    \State $f_x \gets f(x)$
    \State \Return $n, x, f_x, a, b$
    \end{algorithmic}

\subsubsection{Algorithm 3: Optimized Trisection-False Position (Opt.TF)}

\captionof{algorithm}{Pseudocode for the Optimized Trisection-False Position (Opt.TF) algorithm. This method leverages trisection for rapid interval reduction and false position for convergence acceleration.}
    \begin{algorithmic}[1]
    \State \textbf{Input:} Function $f$, interval $[a, b]$, tolerance $tol$, max iterations $max\_iter$
    \State \textbf{Output:} Root approximation $x$, function value $f(x)$, final interval $[a, b]$
    \State $f_a \gets f(a)$, $f_b \gets f(b)$, $\epsilon \gets 10^{-20}$
    \State \textit{// Check if either bound is a root}
    \If{$|f_a| \leq tol$}
        \State \Return $1, a, f_a, a, b$
    \EndIf
    \If{$|f_b| \leq tol$}
        \State \Return $1, b, f_b, a, b$
    \EndIf
    \For{$n \gets 1$ \textbf{to} $max\_iter$}
        \State \textit{// 1. Trisection Step}
        \State $diff \gets b - a$
        \State $x_1 \gets a + diff/3$
        \State $x_2 \gets b - diff/3$
        \State $f_{x_1} \gets f(x_1)$, $f_{x_2} \gets f(x_2)$
        \If{$|f_{x_1}| \leq tol$}
            \State \Return $n, x_1, f_{x_1}, a, b$
        \EndIf
        \If{$|f_{x_2}| \leq tol$}
            \State \Return $n, x_2, f_{x_2}, a, b$
        \EndIf
        
        \State \textit{// Update interval based on trisection}
        \If{$f_a \times f_{x_1} < 0$}
            \State $a, b, f_b \gets a, x_1, f_{x_1}$
        \ElsIf{$f_{x_1} \times f_{x_2} < 0$}
            \State $a, b, f_a, f_b \gets x_1, x_2, f_{x_1}, f_{x_2}$
        \Else
            \State $a, f_a \gets x_2, f_{x_2}$
        \EndIf
    
        \State \textit{// 2. False Position Step with exception handling}
        \State $dx \gets (a \times f_b) - (b \times f_a)$
        \State $dd \gets f_b - f_a$
        \State \textbf{try:}
        \State \quad $x_F \gets dx / dd$
        \State \textbf{catch} (division by zero, overflow, value error):
        \State \quad $x_F \gets dx / (dd + \epsilon)$
        \State $f_{x_F} \gets f(x_F)$
        \If{$|f_{x_F}| \leq tol$}
            \State \Return $n, x_F, f_{x_F}, a, b$
        \EndIf
    
        \State \textit{// Update interval based on false position}
        \If{$f_a \times f_{x_F} < 0$}
            \State $b, f_b \gets x_F, f_{x_F}$
        \Else
            \State $a, f_a \gets x_F, f_{x_F}$
        \EndIf
    \EndFor
    \State \textit{// Max iterations reached}
    \State $x \gets (a+b)/2$
    \State $f_x \gets f(x)$
    \State \Return $max\_iter, x, f_x, a, b$
    \end{algorithmic}

\subsubsection{Algorithm 4: Optimized Trisection-False Position with Modified Secant (Opt.TFMS)}

\captionof{algorithm}{Pseudocode for the Optimized Trisection-False Position with Modified Secant (Opt.TFMS) algorithm. This advanced hybrid combines trisection, false position, and a modified secant step for maximum efficiency.}
    \begin{algorithmic}[1]
    \State \textbf{Input:} Function $f$, interval $[a, b]$, tolerance $tol$, max iterations $max\_iter$
    \State \textbf{Output:} Root approximation $x$, function value $f(x)$, final interval $[a, b]$
    \State $f_a \gets f(a)$, $f_b \gets f(b)$, $\epsilon \gets 10^{-20}$
    \State \textit{// Check if either bound is a root}
    \If{$|f_a| \leq tol$}
        \State \Return $1, a, f_a, a, b$
    \EndIf
    \If{$|f_b| \leq tol$}
        \State \Return $1, b, f_b, a, b$
    \EndIf
    \If{$f_a \times f_b \geq 0$}
        \State \Return $\text{None}, \text{None}, \text{None}, \text{None}, \text{None}$
    \EndIf
    \For{$n \gets 1$ \textbf{to} $max\_iter$}
        \State \textit{// 1. Trisection Step}
        \State $diff \gets b - a$
        \State $x_1 \gets a + diff/3$, $x_2 \gets b - diff/3$
        \State $f_{x_1} \gets f(x_1)$, $f_{x_2} \gets f(x_2)$
        \If{$|f_{x_1}| \leq tol$} \State \Return $n, x_1, f_{x_1}, a, b$ \EndIf
        \If{$|f_{x_2}| \leq tol$} \State \Return $n, x_2, f_{x_2}, a, b$ \EndIf
        
        \State \textit{// Update interval based on trisection results}
        \If{$f_a \times f_{x_1} < 0$}
            \State $b, f_b \gets x_1, f_{x_1}$
        \ElsIf{$f_{x_1} \times f_{x_2} < 0$}
            \State $a, b, f_a, f_b \gets x_1, x_2, f_{x_1}, f_{x_2}$
        \Else
            \State $a, f_a \gets x_2, f_{x_2}$
        \EndIf
    
        \State \textit{// 2. False Position Step with exception handling}
        \State $dx \gets (a \times f_b) - (b \times f_a)$
        \State \textbf{try:}
        \State \quad $x_F \gets dx / (f_b - f_a)$
        \State \textbf{catch} (division by zero, overflow, value error):
        \State \quad $x_F \gets dx / ((f_b - f_a) + \epsilon)$
        \State $f_{x_F} \gets f(x_F)$
        \If{$f_a \times f_{x_F} < 0$}
            \State $b, f_b \gets x_F, f_{x_F}$
        \Else
            \State $a, f_a \gets x_F, f_{x_F}$
        \EndIf
        \If{$|f_{x_F}| \leq tol$}
            \State \Return $n, x_F, f_{x_F}, a, b$
        \EndIf
    
        \State \textit{// 3. Modified Secant Step with exception handling}
        \State $\delta \gets 10^{-8} \max(1, |x_F|) + 10^{-15}$
        \State \textbf{try:}
        \State \quad $x_S \gets x_F - \delta \times f_{x_F} / (f(x_F + \delta) - f_{x_F})$
        \State \textbf{catch} (division by zero, overflow, value error):
        \State \quad $x_S \gets x_F - \delta \times f_{x_F} / ((f(x_F + \delta) - f_{x_F}) + \epsilon)$
        \If{$a < x_S < b$}
            \State $f_{x_S} \gets f(x_S)$
            \If{$|f_{x_S}| < |f_{x_F}|$}
                \If{$f_a \times f_{x_S} < 0$}
                    \State $b, f_b \gets x_S, f_{x_S}$
                \Else
                    \State $a, f_a \gets x_S, f_{x_S}$
                \EndIf
                \If{$|f_{x_S}| \leq tol$}
                    \State \Return $n, x_S, f_{x_S}, a, b$
                \EndIf
            \EndIf
        \EndIf
    \EndFor
    \State \textit{// Max iterations reached}
    \State $x \gets (a+b) \times 0.5$
    \State $f_x \gets f(x)$
    \State \Return $max\_iter, x, f_x, a, b$
    \end{algorithmic}

\setlength{\parskip}{1.5ex plus 0.5ex minus 0.2ex}
\section{Theoretical Analysis}\label{sec:theoretical_analysis}

A primary contribution of this work is a rigorous theoretical analysis of the convergence properties of the proposed hybrid algorithms. While the Intermediate Value Theorem guarantees their eventual convergence, a specific analysis of the multi-phase structure is required to establish the order of convergence and formally characterize the performance bounds.

\subsection{Order of Convergence}

Let $\alpha$ be a simple root of $f(x)$, where $f''$ is continuous in the neighborhood of $\alpha$, and let $e_n = x_n - \alpha$ be the error of an iterate $x_n$. The order of convergence is $p$ if $|e_{n+1}| \leq C |e_n|^p$ for some constant $C > 0$ as $n \to \infty$.

\subsubsection{Analysis of Opt.BF and Opt.TF}

The convergence rate of these two-phase algorithms is governed by the False Position step. For a sequence of iterates generated by False Position, the error recurrence leads to a well-established superlinear order of convergence of \textbf{$p \approx 1.618$}. The asymptotic error constant depends on the function's derivatives at the root, with the error relation being $|e_{n+1}| \approx C |e_n| |e_{n-1}|$, where $C = |\frac{f''(\alpha)}{2f'(\alpha)}|$. Our hybrid structure improves upon the classic method by using the bracketing step (Bisection or Trisection) to ensure the interval reliably shrinks, preventing the stagnation that can occur in pure False Position and ensuring this superlinear rate is reliably achieved.

\subsubsection{Analysis of Opt.BFMS and Opt.TFMS}

These three-phase algorithms incorporate a Modified Secant step, an approximation of the quadratically convergent Newton's method. The error relation for Newton's method is given by $e_{n+1} = C e_n^2 + O(e_n^3)$. Our finite difference approximation of the derivative introduces a perturbation, as shown by its Taylor expansion:
Our finite difference approximation of the derivative introduces a perturbation, as shown by its Taylor expansion:
$$  
    \frac{f(x_n+\delta) - f(x_n)}{\delta} = f'(x_n) + \frac{\delta}{2}f''(x_n) + O(\delta^2)
  $$
The update rule is therefore a perturbed Newton step. With an adaptive perturbation $$ \delta \approx \sqrt{\epsilon_{\text{mach}}} \cdot |x_n| $$, the method maintains near-quadratic convergence while avoiding round-off errors, unlike fixed $$ \delta $$ which can degrade to linear in the terminal phase \cite[e.g.,][]{kelley1995iterative}. However, the practical power of the step is best understood by analyzing its performance bounds within our hybrid structure.

\textbf{Best-Case Performance (Near-Quadratic Acceleration):} When the Modified Secant step is consistently accepted, it becomes the primary driver of convergence. Its error recurrence can be expressed as:
\begin{equation}
    e_{n,S} \approx \frac{f''(\alpha)}{2f'(\alpha)} e_{n,F}^2 - \frac{\delta}{2} \frac{f''(\alpha)}{f'(\alpha)} e_{n,F}
\end{equation}
This relation formally shows that the convergence is driven by a powerful quadratic term but is ultimately limited by a linear error term proportional to $\delta$, resulting in the "effectively quadratic" behavior observed empirically.

\textbf{Worst-Case Performance (Graceful Degradation):} We now provide a brief proof sketch for the algorithm's performance when the Modified Secant step is rejected.

\textit{Proof Sketch.} Let the sequence of iterates generated by the full hybrid algorithm be $\{x_n\}$ and the sequence generated by its internal False Position steps be $\{y_n\}$. The convergence order of $\{y_n\}$ is $p_{FP} \approx 1.618$. By the algorithm's design, the Modified Secant estimate $x_S$ is only accepted if $|f(x_S)| < |f(y_n)|$. If it is rejected, the next state is determined by $y_n$. Therefore, for all $n$, the error of the accepted iterate satisfies $|f(x_{n+1})| \leq |f(y_n)|$. Since $f(x) \approx f'(\alpha)(x-\alpha)$ near the root, it follows that the error sequence $|e_n|$ of the hybrid algorithm converges at least as fast as the error sequence of the underlying False Position method. Thus, the order of convergence $p$ of the hybrid algorithm is bounded below by that of False Position: $p \ge p_{FP} \approx 1.618$. This guarantees at least superlinear convergence.

\textbf{Conclusion on Convergence:} The practical convergence of Opt.BFMS and Opt.TFMS is an adaptive process that opportunistically leverages near-quadratic steps, while the architecture provides a formal guarantee of graceful degradation to a robust superlinear lower bound. This property is a key feature of our multi-phase design.
\subsection{Finite-Precision Considerations and Numerical Stability}

The primary source of instability in these methods is the denominator of the interpolation formulas. For both the false position and modified secant steps, these take the form $f(y) - f(z)$ for two points $y$ and $z$. In standard floating-point arithmetic (IEEE 754 double precision), when $y$ and $z$ are very close to each other, their function values $f(y)$ and $f(z)$ will also be very close, so when perform The subtraction of these two nearly equal numbers, it can lead to \textbf{catastrophic cancellation}, a phenomenon where a large portion of the significant digits are lost, resulting in a denominator dominated by relative error. Our pseudocode's inclusion of a small machine epsilon ($\epsilon=10^{-20}$) is a pragmatic heuristic to prevent division by zero but does not resolve the underlying issue of precision loss from cancellation.

Furthermore, the choice of the perturbation parameter $\delta$ in the modified secant method is critical and represents a fundamental trade-off between both truncation and round-off errors.
\begin{itemize}
    \item \textbf{Truncation Error:} If $\delta$ is too large, the finite difference $\frac{f(x+\delta) - f(x)}{\delta}$ is a poor approximation of the derivative $f'(x)$, and the method's convergence slows. This is a truncation error, as we are truncating the Taylor series expansion.
    \item \textbf{Round-off Error:} If $\delta$ is too small, it becomes comparable to the machine precision. The term $f(x+\delta)$ will be very close to $f(x)$, and their subtraction will suffer from the catastrophic cancellation described above. This is a round-off error.
\end{itemize}

An optimal choice for $$ \delta $$ balances these two error sources, and we use the common heuristic $$ \delta \approx 10^{-8} \cdot \max(1.0, |x|) + 10^{-15} $$, where the scaling approximates $$ \sqrt{\epsilon_{\text{mach}}} \cdot |x| $$ and $$ \epsilon_{\text{mach}} $$ is the machine epsilon (approx. $10^{-16}$ for double precision). This adaptive choice provides robust performance and stability across varying root scales in our benchmark. A candidate $$ x_S $$ is accepted if and only if it resides strictly within the current bracket and yields a smaller function magnitude than $$ x_{curr} $$. The impact of these finite-precision effects is empirically evaluated in Section~\ref{sec:experimental}.

\section{Experimental Setup}\label{sec:experimental}

We designed a rigorous experimental framework to evaluate the proposed algorithms against established baselines (Bisection, False Position) and state-of-the-art methods (Brent's Method, QIR).

\subsection{Benchmark Test Suite}

The algorithms were stress-tested on a diverse suite of \textbf{48 benchmark functions}. This dataset was constructed to cover a wide spectrum of mathematical behaviors, ranging from simple continuous functions to pathological cases with high-frequency oscillations, singularities, and high-degree polynomials (up to degree 25).

The test functions are categorized by their dominant mathematical characteristics:

\textbf{Transcendental Functions:}
\begin{itemize}
\item $f_1(x) = x e^x - 7$: A transcendental function combining exponential and linear terms, interval $[1, 2]$
\item $f_8(x) = e^x - 3x - 2$: A mixed exponential-linear function, interval $[2, 3]$
\item $f_{15}(x) = e^x - 3x^2$: Exponential-quadratic combination, interval $[0, 1]$
\item $f_{18}(x) = x^2 - e^x - 3x + 2$: Quadratic-exponential-linear combination, interval $[0, 1]$
\item $f_{20}(x) = e^{x^2} - e^{\sqrt{2}x}$: Complex exponential function, interval $[0.5, 1.5]$
\item $f_{29}(x) = e^{-x} - \cos(3x) - 0.5$: Mixed exponential-trigonometric function, interval $[0, 1]$
\item $f_{39}(x) = e^{x^2} - e^{1.2x}$: Exponential function with different exponents, interval $[0, 2]$
\item \textbf{Transcendental growth model (Financial Volatility Proxy)}\[ f_46(x) = x e^{x/2} - 1.5, \quad I = [0.0, 1.0] \]
\end{itemize}

\textbf{High-Degree Polynomial Functions:}
\begin{itemize}
\item $f_2(x) = x^3 - x - 1$: A cubic polynomial with multiple roots, interval $[1, 2]$
\item $f_{12}(x) = x^{10} - 1$: A tenth-degree polynomial function, interval $[0, 1.3]$
\item $f_{14}(x) = x^3 - 2x - 5$: Another cubic polynomial, interval $[2, 3]$
\item $f_{19}(x) = (x - 1)^3 + 4(x - 1)^2 - 10$: A cubic polynomial, interval $[0, 3]$
\item $f_{21}(x) = (x^2 - 4)(x + 1.5)(x - 0.5)$: A quartic polynomial, interval $[0, 2]$
\item $f_{22}(x) = x^3 - 3x^2 - 4x + 13$: A cubic polynomial, interval $[-3, -2]$
\item $f_{26}(x) = x^5 - 3x^4 + 25$: A fifth-degree polynomial, interval $[-3, -1]$
\item $f_{27}(x) = x^4 - 2x^2 - 4$: A quartic polynomial, interval $[1, 3]$
\item $f_{32}(x) = x^4 + 2x^2 - x - 1$: A quartic polynomial, interval $[-0.5, 0]$
\item $f_{33}(x) = x^4 - 10x^3 + 35x^2 - 50x + 24$: A quartic polynomial, interval $[0, 1.5]$
\item $f_{35}(x) = x^{25} - 1$: A twenty-fifth-degree polynomial (ill-conditioned), interval $[0, 2]$
\item $f_{38}(x) = x^4 + 2x^3 - 13x^2 - 14x + 24$: A quartic polynomial, interval $[-3, 1]$
\item $f_{40}(x) = x^5 - 3x^4 + 2x^3 - x + 0.1$: A fifth-degree polynomial with small coefficients, interval $[-1, 3]$
\item $f_{41}(x) = (x + 0.3)^7 - 0.01$: A seventh-degree polynomial with small perturbation, interval $[-2, 1]$
\item $f_{42}(x) = x^6 - 8x^5 + 24x^4 - 32x^3 + 16x^2$: A sixth-degree polynomial, interval $[0, 4]$
\item \textbf{Van der Waals equation of state for CO$_2$ (Thermodynamics)} \[ f_45(x) = \left(10 + \frac{3.592}{x^2}\right)(x - 0.04267) - 24.618, \quad I = [2.0, 3.0] \]
\end{itemize}

\textbf{Quadratic Functions:}
\begin{itemize}
\item $f_3(x) = x^2 - x - 2$: A standard quadratic equation, interval $[1, 4]$
\item $f_5(x) = x^2 - 10$: A simple quadratic with irrational roots, interval $[3, 4]$
\item $f_{13}(x) = x^2 + 2x - 7$: A quadratic with different coefficients, interval $[1, 3]$
\item $f_{23}(x) = -0.9x^2 + 1.7x + 2.5$: A quadratic with negative leading coefficient, interval $[2.8, 3.0]$
\item $f_{24}(x) = 1 - 0.61x$: A linear function, interval $[1.5, 2.0]$
\end{itemize}

\textbf{Trigonometric Functions:}
\begin{itemize}
\item $f_4(x) = x - \cos(x)$: A mixed linear-trigonometric function, interval $[0, 1]$
\item $f_6(x) = \sin(x) - x^2$: A trigonometric-quadratic combination, interval $[0.5, 1]$
\item $f_{10}(x) = x \sin(x) - 1$: A product of linear and trigonometric terms, interval $[0, 2]$
\item $f_{11}(x) = x \cos(x) + 1$: A product of linear and trigonometric terms, interval $[-2, 4]$
\item $f_{16}(x) = \sin(10x) - 0.5x$: A high-frequency trigonometric function, interval $[0.1, 0.4]$
\item $f_{17}(x) = x - 0.8\sin(x) - 1.2$: A mixed linear-trigonometric function, interval $[1, 3]$
\item $f_{25}(x) = x^2 |\sin(x)| - 4.1$: A trigonometric function with absolute value, interval $[0, 4]$
\item $f_{28}(x) = x - 0.5\sin(x) - 1$: A mixed linear-trigonometric function, interval $[0, 3]$
\item $f_{34}(x) = 4\sin(x) - x + 1$: A trigonometric-linear combination, interval $[-1, 0]$
\item $f_{37}(x) = \sin(20x) - 0.3x$: A very high-frequency trigonometric function, interval $[0.05, 0.25]$
\item \textbf{Non-linear beam deflection (Structural Engineering)} \[ f_47(x) = \cos(x) + (1 - \cos(x))^2 - 0.05 x^2, \quad I = [3.0, 6.0] \]
\item \textbf{Kepler's equation for high eccentricity (Orbital Mechanics)} \[ f_44(x) = x - 0.99 \sin(x) - 2.0, \quad I = [2.0, 3.0] \]
\end{itemize}

\textbf{Logarithmic and Mixed Functions:}
\begin{itemize}
\item $f_7(x) = x + \ln(x)$: A mixed linear-logarithmic function, interval $[0.1, 1]$
\item $f_9(x) = x^2 + e^{x/2} - 5$: A quadratic-exponential combination, interval $[1, 2]$
\item \textbf{Colebrook-White friction factor (Fluid Mechanics)} \[ f_43(x) = -2 \log_{10}\left(2.7 \times 10^{-5} + \frac{2.51}{10^7 \sqrt{x}}\right) - \frac{1}{\sqrt{x}}, \quad I = [0.008, 0.03] \]
\end{itemize}

\textbf{Ill-Conditioned and Challenging Functions:}
\begin{itemize}
\item $f_{30}(x) = (x - 1)(x - 2)\cdots(x - 20)$: A twentieth-degree polynomial with integer roots, interval $[0, 1.5]$
\item $f_{31}(x) = (x - 1)(x - 2)\cdots(x - 20)$: A twentieth-degree polynomial with integer roots in a different interval, $[19, 21]$
\item $f_{36}(x) = (x - 1.8)^6 (x - 1.81)$: An ill-conditioned polynomial with close roots, interval $[0, 2]$
\item \textbf{Chemical equilibrium reaction extent} \[ f_{48}(x) = 0.05 - \frac{x^3}{(1 - x)(0.8 - 2x)^2}, \quad I = [0.01, 0.3] \]
\end{itemize}

Each function was evaluated over carefully selected intervals that bracket the root, with initial intervals chosen to ensure the Intermediate Value Theorem conditions are satisfied. The expanded dataset includes challenging problems with ill-conditioned roots, high-degree polynomials up to degree 25, and complex transcendental functions, providing a comprehensive and rigorous foundation for algorithm evaluation that addresses the reviewers' concerns about test complexity.

\subsection{Implementation and Compilation Pipeline}

To ensure a fair assessment of computational efficiency, all algorithms were implemented using the SymPy library for symbolic definition but were \textbf{pre-compiled} for execution. 

\textbf{Symbolic Function Definition:} The functions are defined using SymPy symbols, allowing for exact mathematical characterization. This ensures that the test functions are represented without initial numerical approximation errors.

\textbf{Compilation for Performance (sp.lambdify):} A critical feature of our experimental setup is the use of \texttt{sp.lambdify} with the default \textbf{NumPy} backend. While SymPy handles the symbolic definition, the actual execution during benchmarking is performed by \textbf{vectorized NumPy} functions generated from the symbolic expressions. This eliminates virtually all interpretation overhead of the symbolic engine itself, so the reported CPU times predominantly reflect the efficiency of the numerical algorithm rather than SymPy’s expression evaluation.

\textbf{Comprehensive Error Handling:} The implementation features extensive error surveillances for rare cases like division by zero, overflow conditions, and convergence failures.

\textbf{Adaptive Parameters:} The modified secant step uses an adaptive $$ \delta $$ computed per iteration for enhanced numerical robustness.

\subsection{Performance Evaluation and metrics}

We employed a strict timing protocol to ensure statistical validity and minimize the impact of system noise (such as OS interrupts or memory management):

\textbf{Single-Loop Strategy:} We utilized a high-volume single-loop strategy. Each algorithm was executed $N=10,000$ times per problem.

\textbf{Garbage Collection Control:} To prevent Python's automatic memory management from skewing microsecond-level measurements, the Garbage Collector (GC) was explicitly disabled immediately prior to the critical execution block and re-enabled immediately after.

\subsection{Metrics:}

\textbf{Number of Iterations}: The number of algorithmic loops required to converge to the root; with fewer iterations generally preferring an algorithm, but this has to be traded for the computation cost per iteration.

\textbf{CPU Times}: These measures the actual computational efficiency. CPU time remains the most viable measure when referring to applications in the real world where resource consumption matters, The execution time was measured using the high-resolution \texttt{time.perf\_counter\_ns()} and converted to microseconds.

\textbf{Function Value at Computed Root:} The accuracy of the root approximation is basically measured by this metric - how much close the computed value is to a true root of the function. A number very close to zero confirms maximum accuracy. 

\textbf{Final Interval Size:} This metric measures how precise the root approximation is, with respect to the actual size of the interval where the computed root lies in. Smaller intervals indicate higher precision.

Convergence tolerance was kept at $10^{-14}$ to ensure very precise results, and each algorithm was terminated after reaching the limit of 10,000 iterations to avoid loop problems within the algorithm. These parameters were chosen according to the features and needs of high-precision numerical analysis as well as the capabilities of the SymPy library.

This methodology ensures that the performance differences reported in Section \ref{sec:results} are attributable to the algorithmic logic itself rather than implementation artifacts.

\subsection{Algorithm Code Legend}
\begin{table}[H]
    \centering
    \captionof{table}{Algorithm code legend used throughout all tables. Each code corresponds to a specific root-finding algorithm, as defined in the methodology.}
        \label{tab:algo_legend}
        \resizebox{\textwidth}{!}{
        \begin{tabular}{ll}
        \hline
        Code & Algorithm Name \\
        \hline
        Bisection & Bisection \\
        False Position & Regula Falsi (False Position) \\
        Trisection & Trisection \\
        HybridBlendTF & Badr-2021 - Hybrid Blend Trisection False Position \\
        HybridBlendBF & Sabharwal 2019 - Blended Root Finding Algorithm Outperforms Bisection and Regula Falsi \\
        Brent (SciPy) & Brent (SciPy implementation) \\
        Brent (Impl.) & Brent (my implementation) \\
        Opt.BF & Optimized Hybrid Bisection False Position \\
        Opt.BFMS & Optimized Hybrid Bisection False Position Modified Secant \\
        Opt.TF & Optimized Hybrid Trisection False Position \\
        Opt.TFMS & Optimized Hybrid Trisection False Position Modified Secant \\
        QIR & Quadratic Interval Refinement (Abbott 2012) \\
        \hline
    \end{tabular}}
    \end{table}

\subsection*{Case Study 1: High-Degree Polynomial (Problem 42)}
To demonstrate performance on high-degree polynomials, we analyze Problem 42 from our test suite. This problem involves the sixth-degree polynomial, a class of functions common in engineering optimization and control systems analysis. The function is:
\[
f(x) = x^6 - 8x^5 + 24x^4 - 32x^3 + 16x^2
\]
on the interval $[0, 4]$. Table~\ref{tab:realworld1} summarizes the results for all algorithms on this problem.

\begin{table}[H]
\centering
\captionof{table}{Performance on the high-degree polynomial $f(x) = x^6 - 8x^5 + 24x^4 - 32x^3 + 16x^2$ over $[0,4]$. Metrics reported are average CPU time (in microseconds) over 10,000 runs.}
\label{tab:realworld1}
\begin{tabular}{lcc}
\hline
Algorithm & CPU Time ($\mu s$) & Relative Performance \\
\hline
Opt.BFMS & 0.8867 & 1.00x (fastest) \\
Opt.BF & 0.8869 & 1.00x \\
Opt.TFMS & 0.8924 & 1.01x \\
Opt.TF & 0.9294 & 1.05x \\
Bisection & 1.0253 & 1.16x \\
QIR & 1.2065 & 1.36x \\
False Position & 1.2323 & 1.39x \\
Brent (Impl.) & 1.2881 & 1.45x \\
HybridBlendBF & 1.7885 & 2.02x \\
HybridBlendTF & 2.2088 & 2.49x \\
Brent (SciPy) & 6.4158 & 7.24x \\
Trisection & 13093.77 & 14766.05x \\
\hline
\end{tabular}
\end{table}

On this high-degree polynomial, the proposed Opt.BFMS algorithm achieves the fastest performance. The results show the effectiveness of the hybrid approach.

\subsection*{Case Study 2: Hardest Benchmark Problem (Problem 30, the classic Wilkinson's polynomial)}
To demonstrate practical applicability on computationally demanding tasks, we analyze Problem 30, the classic Wilkinson's polynomial, the most challenging function in our test suite. This problem is a twentieth-degree polynomial with known integer roots:
\[
f(x) = \prod_{k=1}^{20} (x - k)
\]
on the interval $[0, 1.5]$. Table~\ref{tab:realworld2} summarizes the results for this worst-case scenario.

\begin{table}[H]
\centering
\captionof{table}{Performance on the hardest benchmark problem, a twentieth-degree polynomial, over $[0,1.5]$. Metrics reported are average CPU time (in microseconds) over 10,000 runs.}
\label{tab:realworld2}
\begin{tabular}{lcc}
\hline
Algorithm & CPU Time ($\mu s$) & Relative Performance \\
\hline
Trisection & 2.1215 & 1.00x (fastest) \\
Opt.TF & 2.9798 & 1.40x \\
Opt.TFMS & 2.9959 & 1.41x \\
HybridBlendTF & 3.3490 & 1.58x \\
Opt.BFMS & 14.1943 & 6.69x \\
Opt.BF & 16.5449 & 7.80x \\
Brent (SciPy) & 24.8929 & 11.73x \\
Brent (Impl.) & 27.0669 & 12.76x \\
QIR & 27.9347 & 13.17x \\
HybridBlendBF & 38.4100 & 18.11x \\
Bisection & 64.1714 & 30.25x \\
False Position & 17905.05 & 8439.66x \\
\hline
\end{tabular}
\end{table}

On this exceptionally difficult problem, the simple Trisection method finds the root most efficiently. However, the proposed trisection-based hybrids, Opt.TF and Opt.TFMS, perform nearly as well, ranking second and third. This demonstrates their value as robust alternatives.

\subsection*{Case Study 3: Cubic Polynomial Test Case (Problem 19)}
To evaluate performance on a standard polynomial function, we analyze Problem 19 from the benchmark set. This case study demonstrates the algorithms' efficiency on a common class of problems found in numerical analysis. The test function is the cubic polynomial:
\[
f(x) = (x - 1)^3 + 4(x - 1)^2 - 10,
\]
on the interval $[0, 3]$. Table~\ref{tab:realworld3} summarizes the results for all algorithms on this problem.

\begin{table}[H]
\centering
\captionof{table}{Performance on the cubic polynomial problem $(x-1)^3 + 4(x-1)^2 - 10$ over $[0,3]$. Metrics reported are average CPU time (in microseconds) over 10,000 runs.}
\label{tab:realworld3}
\begin{tabular}{lcc}
\hline
Algorithm & CPU Time ($\mu s$) & Relative Performance \\
\hline
Opt.BFMS & 5.0414 & 1.00x (fastest) \\
Opt.TFMS & 5.5388 & 1.10x \\
Opt.BF & 6.4508 & 1.28x \\
Opt.TF & 7.2704 & 1.44x \\
Brent (Impl.) & 12.3514 & 2.45x \\
QIR & 13.1685 & 2.61x \\
HybridBlendBF & 14.7042 & 2.92x \\
HybridBlendTF & 14.7985 & 2.94x \\
Brent (SciPy) & 16.2649 & 3.23x \\
False Position & 20.4811 & 4.06x \\
Bisection & 26.7110 & 5.30x \\
Trisection & 27.5942 & 5.47x \\
\hline
\end{tabular}
\end{table}

The results again highlight the efficiency of the optimized hybrid methods, with Opt.BFMS leading the performance on this cubic polynomial, followed closely by the other proposed hybrid variants. 
    
\subsection*{Case Study 4: Kepler's Equation (Problem 44)}
To further demonstrate algorithm behaviour on a representative transcendental test from Orbital Mechanics, we consider Problem 44 from the benchmark set:
\[
f(x) = x - 0.99 \sin(x) - 2.0,
\]
on the interval $[2.0, 3.0]$. Table~\ref{tab:realworld4} summarizes the average CPU time (in microseconds) for each algorithm on this problem and their relative performance compared to the fastest method.

\begin{table}[H]
\centering
\captionof{table}{Performance on $f(x) = x - 0.99 \sin(x) - 2.0$ over the interval $[2.0, 3.0]$. Metrics reported are average CPU time (in $\mu s$) over 10,000 runs.}
\label{tab:realworld4}
\begin{tabular}{lcc}
\hline
Algorithm & CPU Time ($\mu s$) & Relative Performance \\
\hline
Opt.BFMS & 10.4138 & 1.00x (fastest) \\
Opt.TFMS & 16.4028 & 1.58x \\
Opt.BF & 17.1239 & 1.64x \\
Opt.TF & 17.5231 & 1.68x \\
Brent (SciPy) & 17.6197 & 1.69x \\
Brent (Impl.) & 25.1547 & 2.42x \\
QIR & 28.5341 & 2.74x \\
False Position & 29.8755 & 2.87x \\
HybridBlendTF & 29.9091 & 2.87x \\
HybridBlendBF & 32.7019 & 3.14x \\
Trisection & 76.9222 & 7.39x \\
Bisection & 84.0384 & 8.07x \\
\hline
\end{tabular}
\end{table}

The results for Problem 44 indicate that the optimized BFMS variant remains the most efficient solver for this high-eccentricity Kepler problem. It significantly outperforms the other methods.

\subsection*{Case Study 5: Chemical Equilibrium (Problem 48)}
To demonstrate performance on ill-conditioned rational functions, we analyze Problem 48 from our test suite. This function describes the extent of reaction in chemical equilibrium and presents a challenge due to the singularity near the boundary.
\[
f_{48}(x) = 0.05 - \frac{x^3}{(1 - x)(0.8 - 2x)^2}
\]
on the interval $[0.01, 0.3]$. Table~\ref{tab:realworld5} summarizes the results for all algorithms on this problem.

\begin{table}[H]
\centering
\captionof{table}{Performance on the chemical equilibrium function $f_{48}(x)$ over $[0.01, 0.3]$. Metrics reported are average CPU time (in microseconds) over 10,000 runs.}
\label{tab:realworld5}
\begin{tabular}{lcc}
\hline
Algorithm & CPU Time ($\mu s$) & Relative Performance \\
\hline
Opt.BFMS & 5.4643 & 1.00x (fastest) \\
Opt.BF & 6.3986 & 1.17x \\
Opt.TF & 6.4557 & 1.18x \\
Opt.TFMS & 6.5797 & 1.20x \\
HybridBlendBF & 12.1746 & 2.23x \\
HybridBlendTF & 13.4335 & 2.46x \\
QIR & 15.7540 & 2.88x \\
Brent (Impl.) & 16.7691 & 3.07x \\
Brent (SciPy) & 18.5895 & 3.40x \\
Bisection & 22.6479 & 4.14x \\
Trisection & 22.8773 & 4.19x \\
False Position & 134.9078 & 24.69x \\
\hline
\end{tabular}
\end{table}

On this challenging rational function, the proposed Opt.BFMS algorithm demonstrates great performance, achieving the fastest convergence.
\FloatBarrier

\section{Results and Analysis}\label{sec:results}

This section presents a comprehensive analysis of the performance data, including detailed comparisons across all algorithms and test functions, statistical significance of the results, and practical implications for numerical computing applications.

\subsection{Overall Performance Analysis}
The experimental results reveal significant performance improvements achieved by the hybrid algorithms. Table \ref{tab:performance_comparison_main} presents the comprehensive performance metrics for all twelve algorithms, including the four proposed hybrid methods, five traditional/baseline methods, and three state-of-the-art methods for comparison.

\begin{table}
\centering
\captionof{table}{Comprehensive comparison of all root-finding algorithms across 48 diverse benchmark problems. Reported metrics are the average CPU time (in microseconds), median CPU time, standard deviation, and number of wins (problems where the algorithm achieved the fastest time). Algorithm codes are defined in Table~\ref{tab:algo_legend}.}
\label{tab:performance_comparison_main}
\resizebox{\textwidth}{!}{
\begin{tabular}{lcccccc}
\hline
Algorithm & Avg CPU Time (us) &  Min CPU Time (us) &  Max CPU Time (us) & Median CPU Time (us) & Std Dev (us) & Wins \\
\hline
Bisection & 61.64110822 & 0.41 & 2065.696 & 61.9375 & 63.65484592 & 2 \\
False Position & 757.4360333 & 0.448 & 35476 & 32.627 & 2791.630865 & 1 \\
Trisection & 333.4573176 & 0.642 & 26974.878 & 69.7175 & 1863.828474 & 3 \\
HybridBlendTF & 24.50749655 & 0.806 & 258.906 & 14.482 & 26.13210366 & 0 \\
HybridBlendBF & 24.86951824 & 0.628 & 260.014 & 15.318 & 27.06712258 & 0 \\
QIR & 25.24653399 & 0.741 & 723.458 & 15.288 & 29.9451163 & 1 \\
Brent (Impl.) & 18.09817806 & 0.58 & 509.633 & 13.712 & 14.77488367 & 3 \\
Brent (SciPy) & 17.89097685 & 5.376 & 560.432 & 17.001 & 8.838381473 & 1 \\
Opt.BF & 12.64507349 & 0.463 & 121.84 & 8.383 & 12.65903175 & 3 \\
Opt.BFMS & 10.25237277 & 0.472 & 151.102 & 5.869 & 10.86245175 & 32 \\
Opt.TF & 15.01685269 & 0.461 & 607.776 & 7.406 & 15.97002554 & 2 \\
Opt.TFMS & 12.44601301 & 0.467 & 631.768 & 6.3575 & 13.68494785 & 0 \\
\hline
\end{tabular}}
\end{table}
\FloatBarrier

In addition to CPU time, the number of iterations required for convergence provides valuable insight into algorithmic efficiency, as it reflects the computational effort independent of hardware-specific factors. Table~\ref{tab:iteration_count_comparison} presents the average iteration count, maximum iteration count , minimum iteration count, median iteration count, standard deviation, and number of wins (problems where an algorithm achieved the lowest iteration count) for each algorithm across the 48 benchmark problems.
\begin{table}
\centering
\captionof{table}{Comprehensive comparison of all root-finding algorithms across 48 diverse benchmark problems. Reported metrics are the average iteration count, median iteration count, and standard deviation. Algorithm codes are defined in Table~\ref{tab:algo_legend}.}
\label{tab:iteration_count_comparison}
\resizebox{\textwidth}{!}{
\begin{tabular}{lcccccc}
\hline
Algorithm & Avg Iteration Count & Min Iteration Count & Max Iteration Count & Median Iteration Count & Std Dev & Wins \\
\hline
Bisection & 41.5625 & 1 & 52 & 47 & 15.14089071 & 3 \\
False Position & 938.8125 & 1 & 10000 & 18 & 2784.503856 & 4 \\
Trisection & 235.0833333 & 1 & 10000 & 29 & 1424.383072 & 4 \\
HybridBlendTF & 5.895833333 & 1 & 10 & 7 & 2.686695184 & 0 \\
HybridBlendBF & 7.125 & 1 & 14 & 8 & 3.269974099 & 0 \\
QIR & 6.375 & 1 & 31 & 5 & 5.897127616 & 3 \\
Brent (Impl.) & 1052.229167 & 1 & 10000 & 13 & 3051.171884 & 0 \\
Brent (SciPy) & 8.979166667 & 1 & 33 & 8.5 & 4.888297156 & 0 \\
Opt.BF & 6.6875 & 1 & 12 & 7 & 2.895086852 & 0 \\
Opt.BFMS & 3.020833333 & 1 & 6 & 3 & 1.145455676 & 27 \\
Opt.TF & 5.291666667 & 1 & 9 & 6 & 2.362542878 & 0 \\
Opt.TFMS & 2.8125 & 1 & 6 & 3 & 1.148482984 & 7 \\
\hline
\end{tabular}}
\end{table}
\FloatBarrier

\begin{table}[H]
    \centering
    \captionof{table}{Comprehensive Number of Function Evaluations (NFE) comparison across 48 problems. Lower NFE indicates higher algorithmic efficiency. Averages for some methods are heavily skewed by non-convergence on certain problems.}
    \label{tab:nfe_comparison_main}
    \begin{tabular}{lcc}
    \hline
    Algorithm & Avg NFE (AVG number of Iterations $\times$ NFE/iter) & Median NFE \\
    \hline
    \textbf{Opt.BFMS} & \textbf{9.062499999} (3.020833333 $\times$ 3) & \textbf{9} \\
    \textbf{Opt.TFMS} & \textbf{11.25} (2.8125 $\times$ 4) & \textbf{12} \\
    Opt.BF & 13.375 (6.6875 $\times$ 2) & 14 \\
    Opt.TF & 15.875 (5.291666667 $\times$ 3) & 18 \\
    \hline
    \end{tabular}
\end{table}

\subsection{Detailed Performance by Problem}

Table \ref{tab:detailed_results} provides a detailed breakdown of CPU times for each algorithm across all test problems, revealing the consistency and robustness of the proposed methods.

\begin{table}[H]
\centering
\captionof{table}{Average CPU time in microseconds (us) for each algorithm on all 48 benchmark problems. Each cell shows the mean CPU time over 10000 runs for the corresponding problem and algorithm. See Table~\ref{tab:algo_legend} for algorithm codes.}
\label{tab:detailed_results}
\resizebox{\textwidth}{!}{
\begin{tabular}{lcccccccccccc}
\hline
Problem & Bisection & False Position & Trisection & HybridBlendTF & HybridBlendBF & QIR & Brent (Impl.) & Brent (SciPy) & Opt.BF & Opt.BFMS & Opt.TF & Opt.TFMS \\
\hline
Problem 1 & 82.03971 & 75.120658 & 88.6167844 & 34.9211277 & 40.6539372 & 17.272257 & 31.6720291 & 19.9974675 & 21.2822234 & 12.4861186 & 23.8169781 & 18.6548579 \\
Problem 2 & 15.1799397 & 16.4431041 & 15.2112764 & 7.9725063 & 7.3770947 & 9.5173385 & 10.6627914 & 13.6142813 & 3.8329213 & 2.8706915 & 3.6736849 & 3.298393 \\
Problem 3 & 15.7291882 & 15.617691 & 0.6844225 & 1.0365921 & 1.8456664 & 3.0621499 & 11.2649888 & 13.593864 & 3.7000002 & 2.8323308 & 0.9709366 & 0.9931395 \\
Problem 4 & 69.2231644 & 26.8285466 & 71.4718546 & 31.4734951 & 29.5029597 & 30.0567027 & 23.7651969 & 16.8716641 & 15.6294641 & 11.270065 & 18.7534461 & 14.9294715 \\
Problem 5 & 13.8024563 & 6.6314817 & 15.9294301 & 7.6313217 & 7.1242153 & 3.0796584 & 8.4712762 & 11.5384367 & 3.3201566 & 2.3365714 & 4.1435888 & 3.2652504 \\
Problem 6 & 78.7983818 & 42.2516652 & 78.3722196 & 25.4427832 & 29.784928 & 13.2405157 & 13.9016562 & 17.5232889 & 19.6613027 & 14.8324417 & 24.021916 & 16.953836 \\
Problem 7 & 77.4173089 & 93.2458618 & 73.1770686 & 34.1106758 & 27.22368 & 17.4306602 & 16.5776441 & 17.571427 & 14.2817129 & 11.7478362 & 20.1988043 & 16.2535763 \\
Problem 8 & 81.5716862 & 124.7168885 & 78.626117 & 36.3101582 & 39.0481729 & 17.5508276 & 39.5168049 & 18.4728279 & 22.325765 & 15.094676 & 24.7939972 & 19.6068608 \\
Problem 9 & 94.848981 & 44.7712509 & 77.7713454 & 33.4528442 & 36.8067168 & 13.9189014 & 26.0553473 & 18.3422752 & 21.2522477 & 13.7575211 & 19.6096021 & 18.5370236 \\
Problem 10 & 80.6010016 & 15.6624168 & 72.943363 & 24.3166514 & 23.5351186 & 63.2375051 & 22.1186477 & 17.2867728 & 12.5225269 & 12.0455225 & 17.4451102 & 15.8207237 \\
Problem 11 & 78.6048521 & 31.0729053 & 73.8380013 & 38.481575 & 39.944846 & 54.0428891 & 18.18286 & 22.3388371 & 18.8209877 & 14.0758581 & 20.1829507 & 18.0321834 \\
Problem 12 & 14.5739218 & 55.4028502 & 16.0456766 & 9.9826461 & 10.6254487 & 13.6407164 & 10.3693458 & 13.5185488 & 4.7746973 & 2.834563 & 5.3116854 & 4.9650616 \\
Problem 13 & 17.0673132 & 10.024112 & 17.7803717 & 8.7960764 & 5.2460148 & 3.1939393 & 8.8884809 & 13.0384864 & 4.1314531 & 2.7475831 & 5.5040751 & 3.7171769 \\
Problem 14 & 17.1326425 & 16.3859693 & 18.3931141 & 9.9143237 & 9.0552154 & 10.0717003 & 10.9997442 & 12.0753094 & 4.6801299 & 3.1190414 & 5.3849582 & 3.7018317 \\
Problem 15 & 82.4596066 & 33.9687556 & 77.8564457 & 32.0340556 & 35.3937286 & 40.2333455 & 14.3820353 & 18.2137961 & 18.2364726 & 15.3324184 & 22.3490211 & 17.6285665 \\
Problem 16 & 68.9317951 & 16.4251713 & 75.4587098 & 25.2577711 & 25.1787312 & 67.2607561 & 13.5432118 & 20.1992457 & 12.9681977 & 12.6205718 & 15.2687737 & 13.4845946 \\
Problem 17 & 77.0955486 & 53.7090668 & 82.3765433 & 39.7518736 & 28.6271591 & 39.4892603 & 31.5412598 & 21.0468991 & 14.9730571 & 12.4553831 & 23.6864182 & 16.3091102 \\
Problem 18 & 90.8065411 & 30.8263643 & 92.2717414 & 34.727939 & 33.3911253 & 18.9936934 & 14.7736222 & 17.0648278 & 17.0686647 & 14.1708554 & 24.2591519 & 14.9777697 \\
Problem 19 & 26.7109945 & 20.4810923 & 27.5941655 & 14.7985002 & 14.7041797 & 13.168528 & 12.3513977 & 16.2648816 & 6.4507984 & 5.0413895 & 7.2703838 & 5.538753 \\
Problem 20 & 235.2006721 & 176.8142546 & 235.893905 & 90.9293856 & 115.0985569 & 41.7632943 & 33.7511921 & 43.286472 & 52.5869682 & 44.5614295 & 60.3051059 & 54.8251388 \\
Problem 21 & 0.9609488 & 0.6661714 & 18.8428798 & 1.1976756 & 0.9360953 & 0.8202185 & 0.6304071 & 5.7221202 & 0.5065221 & 0.5092239 & 0.5114453 & 0.5030884 \\
Problem 22 & 26.9972603 & 21.6333657 & 27.9174218 & 13.5069334 & 12.2793174 & 13.6895045 & 11.6793859 & 13.0345745 & 6.7594997 & 4.5004427 & 7.5696629 & 5.0690929 \\
Problem 23 & 14.2964057 & 4.6642207 & 15.7906737 & 5.7105611 & 5.6398458 & 3.1473387 & 6.6206719 & 10.5852162 & 3.0387762 & 2.1607546 & 3.1866073 & 2.4176974 \\
Problem 24 & 9.9416251 & 0.4917175 & 12.0964428 & 0.8630455 & 0.674448 & 1.930598 & 0.9222016 & 6.3879042 & 0.9147416 & 0.7671825 & 1.0803668 & 1.1162384 \\
Problem 25 & 150.9066032 & 155.2129174 & 146.2081248 & 65.1462125 & 66.573632 & 42.6410517 & 51.5820465 & 32.6008789 & 33.9379473 & 30.23272 & 38.2469087 & 33.0184345 \\
Problem 26 & 24.2208063 & 6558.77945 & 26.5286644 & 12.5400864 & 11.1753848 & 15.8596035 & 13.2173431 & 17.0889747 & 5.7848169 & 4.4985205 & 7.0120083 & 5.9046661 \\
Problem 27 & 22.1768397 & 53.4443911 & 23.2692108 & 12.8421894 & 10.5857542 & 12.6050298 & 12.5045425 & 15.7632703 & 5.5120794 & 4.6568877 & 7.0888683 & 4.6788031 \\
Problem 28 & 80.5009844 & 64.8425199 & 82.7027534 & 34.9555255 & 20.4632311 & 13.4512636 & 31.2486925 & 19.9813091 & 10.7108641 & 10.3733646 & 23.8169004 & 18.7384068 \\
Problem 29 & 147.8446401 & 44.0434802 & 138.9692749 & 61.7432812 & 49.6894514 & 100.2587049 & 43.0690606 & 28.8537162 & 25.7923673 & 21.7610216 & 35.9287243 & 28.695248 \\
Problem 30 & 64.1714252 & 17905.04867 & 2.1215363 & 3.3490099 & 38.4099573 & 27.9347408 & 27.0668598 & 24.8928985 & 16.5448848 & 14.1942997 & 2.9798294 & 2.9959257 \\
Problem 31 & 1.4200666 & 1.9960233 & 60.4178539 & 3.3301583 & 2.6578137 & 1.6606999 & 1.9998479 & 7.0054936 & 1.3540572 & 1.3950537 & 1.3944482 & 1.3652428 \\
Problem 32 & 22.0114611 & 6.387663 & 22.9065344 & 9.7889377 & 7.9728705 & 11.1881492 & 3.8578385 & 11.6487834 & 4.3691065 & 3.6785412 & 5.3310753 & 4.6917274 \\
Problem 33 & 29.9191323 & 98.9769685 & 1.0803418 & 1.7734505 & 15.5445114 & 14.5711609 & 14.5504146 & 17.0884208 & 8.233629 & 5.3305733 & 1.5576255 & 1.7128804 \\
Problem 34 & 84.7105786 & 33.0761412 & 82.0800047 & 30.9821794 & 34.1113266 & 39.7060014 & 29.9458659 & 18.0812022 & 17.8932824 & 12.9280055 & 18.4076703 & 13.9959406 \\
Problem 35 & 0.4466111 & 4078.921644 & 17.3851043 & 11.3160822 & 0.7884049 & 1.1576122 & 2.5120656 & 10.5124278 & 0.7941358 & 0.8688168 & 6.1317199 & 6.0771559 \\
Problem 36 & 1.8056138 & 4270.481172 & 2.359485 & 4.6199349 & 4.7207693 & 14.9515541 & 10.9896909 & 36.5244152 & 2.9530384 & 4.7177436 & 2.9008822 & 4.0028582 \\
Problem 37 & 77.2001783 & 16.3904413 & 77.7461541 & 25.2321979 & 29.6408043 & 70.2874548 & 16.1554465 & 19.1719755 & 12.9692218 & 12.5857559 & 14.9631398 & 16.8094476 \\
Problem 38 & 1.5360861 & 1.0138045 & 34.8969022 & 1.8566177 & 1.4153531 & 1.0629029 & 1.0458473 & 6.1464935 & 0.7752547 & 0.7975989 & 0.8007951 & 0.7878999 \\
Problem 39 & 151.2999822 & 4.9771934 & 149.8032217 & 8.2268323 & 6.7786029 & 3.6260888 & 4.6496416 & 9.794647 & 3.1789701 & 3.2275837 & 3.1635904 & 3.3498438 \\
Problem 40 & 28.2289317 & 102.7544173 & 28.9265376 & 15.8196775 & 13.9550418 & 18.1758722 & 18.969367 & 16.9835311 & 8.8255142 & 6.0026698 & 9.7880467 & 7.2097865 \\
Problem 41 & 12.8526571 & 1532.498148 & 14.5110377 & 9.8109086 & 10.6156944 & 15.9829847 & 9.6112404 & 21.3617878 & 4.0805906 & 3.6960205 & 5.2180281 & 5.0395907 \\
Problem 42 & 1.0253341 & 1.2322899 & 13093.76598 & 2.2087552 & 1.7884532 & 1.2065451 & 1.2880809 & 6.4158027 & 0.8869292 & 0.8867468 & 0.9293502 & 0.8923835 \\
Problem 43 & 320.3543609 & 254.8126202 & 306.7669142 & 135.6978444 & 138.3829517 & 51.3516863 & 71.5150089 & 43.6915765 & 61.9720426 & 55.0492306 & 79.7639155 & 70.8835261 \\
Problem 44 & 84.0383725 & 29.8755134 & 76.9221696 & 29.909085 & 32.7019178 & 28.5340548 & 25.1546966 & 17.6197232 & 17.1238741 & 10.4138229 & 17.5230918 & 16.402799 \\
Problem 45 & 15.7712929 & 4.0096286 & 18.7544441 & 7.45702 & 5.917869 & 13.9315391 & 6.7946844 & 10.7295404 & 2.8901668 & 2.2920424 & 4.1741079 & 2.621474 \\
Problem 46 & 85.2255146 & 27.9673941 & 77.5340715 & 27.3363024 & 29.4922612 & 17.2070567 & 11.8118418 & 16.0811602 & 15.8454762 & 12.9740213 & 22.6358324 & 17.1375105 \\
Problem 47 & 158.4658745 & 41.4237377 & 132.4576014 & 84.3635439 & 68.4829805 & 168.9155605 & 49.7610923 & 34.549938 & 34.4174258 & 29.9180535 & 45.2979644 & 33.2179499 \\
Problem 48 & 22.6479016 & 134.9077875 & 22.8773458 & 13.433484 & 12.1746367 & 15.7540149 & 16.7691319 & 18.5894977 & 6.398565 & 5.4643265 & 6.4557354 & 6.5796875 \\
\hline
\end{tabular}}
\end{table}
\FloatBarrier

Table~\ref{tab:iteration_count_results} presents the exact number of iterations for each algorithm on each problem, revealing patterns of convergence behavior. 

\noindent\textbf{Note on non-convergence:} An iteration count of \textbf{10000} indicates that the algorithm \textbf{did not converge to the specified tolerance} ($10^{-14}$) and instead reached the \textbf{maximum iteration cap} before approaching the tolerance.

\begin{table}
\centering
\captionof{table}{Iteration count for each algorithm on all 48 benchmark problems. Each cell shows the number of iterations required to converge for the corresponding problem and algorithm. See Table~\ref{tab:algo_legend} for algorithm codes.  Brent (SciPy) is excluded due to unavailable iteration count data.}
\label{tab:iteration_count_results}
\resizebox{\textwidth}{!}{
\begin{tabular}{lcccccccccccc}
\hline
Prpoblem & Bisection & False Position & Trisection & HybridBlendTF & HybridBlendBF & QIR & Brent (Impl.) & Brent (SciPy) & Opt.BF & Opt.BFMS & Opt.TF & Opt.TFMS  \\
\hline
Problem 1 & 45 & 28 & 32 & 7 & 10 & 4 & 13 & 9 & 9 & 3 & 7 & 3  \\
Problem 2 & 48 & 39 & 28 & 7 & 8 & 4 & 15 & 10 & 8 & 3 & 5 & 3  \\
Problem 3 & 49 & 37 & 1 & 1 & 2 & 1 & 16 & 10 & 8 & 3 & 1 & 1  \\
Problem 4 & 44 & 11 & 29 & 7 & 8 & 8 & 10 & 8 & 7 & 3 & 6 & 3  \\
Problem 5 & 47 & 16 & 31 & 7 & 8 & 1 & 12 & 8 & 7 & 3 & 6 & 3  \\
Problem 6 & 45 & 16 & 29 & 5 & 7 & 3 & 5 & 8 & 8 & 3 & 7 & 3  \\
Problem 7 & 47 & 37 & 28 & 7 & 7 & 4 & 7 & 8 & 6 & 3 & 6 & 3  \\
Problem 8 & 46 & 44 & 28 & 7 & 9 & 4 & 16 & 8 & 9 & 3 & 7 & 3  \\
Problem 9 & 48 & 15 & 26 & 6 & 8 & 3 & 10 & 8 & 8 & 3 & 5 & 3  \\
Problem 10 & 46 & 6 & 28 & 5 & 6 & 17 & 9 & 8 & 5 & 3 & 5 & 3  \\
Problem 11 & 45 & 12 & 28 & 8 & 10 & 14 & 7 & 11 & 8 & 3 & 6 & 3  \\
Problem 12 & 50 & 138 & 31 & 9 & 12 & 6 & 15 & 10 & 11 & 4 & 8 & 5  \\
Problem 13 & 48 & 20 & 29 & 7 & 5 & 1 & 12 & 9 & 7 & 3 & 7 & 3  \\
Problem 14 & 49 & 33 & 30 & 8 & 9 & 4 & 15 & 8 & 9 & 3 & 7 & 3  \\
Problem 15 & 46 & 12 & 27 & 6 & 8 & 10 & 5 & 8 & 7 & 3 & 6 & 3  \\
Problem 16 & 42 & 6 & 29 & 5 & 6 & 18 & 5 & 10 & 5 & 3 & 4 & 2  \\
Problem 17 & 44 & 20 & 30 & 8 & 7 & 10 & 13 & 10 & 6 & 3 & 7 & 3  \\
Problem 18 & 47 & 10 & 30 & 6 & 7 & 4 & 5 & 7 & 6 & 3 & 6 & 2  \\
Problem 19 & 51 & 28 & 32 & 9 & 11 & 5 & 15 & 11 & 9 & 4 & 7 & 3  \\
Problem 20 & 48 & 24 & 31 & 7 & 11 & 5 & 7 & 11 & 9 & 4 & 7 & 4  \\
Problem 21 & 2 & 1 & 31 & 1 & 1 & 1 & 10000 & 1 & 1 & 1 & 1 & 1  \\
Problem 22 & 51 & 29 & 32 & 8 & 9 & 5 & 15 & 8 & 10 & 3 & 7 & 3  \\
Problem 23 & 45 & 10 & 28 & 5 & 6 & 1 & 9 & 7 & 6 & 2 & 4 & 2  \\
Problem 24 & 43 & 1 & 28 & 1 & 1 & 1 & 1 & 3 & 1 & 1 & 1 & 1  \\
Problem 25 & 51 & 34 & 32 & 8 & 10 & 8 & 17 & 12 & 9 & 4 & 7 & 4  \\
Problem 26 & 51 & 10000 & 34 & 8 & 9 & 6 & 18 & 12 & 9 & 4 & 7 & 4  \\
Problem 27 & 51 & 90 & 32 & 9 & 9 & 5 & 16 & 11 & 9 & 4 & 8 & 3  \\
Problem 28 & 46 & 24 & 30 & 7 & 5 & 3 & 13 & 9 & 4 & 2 & 7 & 3  \\
Problem 29 & 47 & 9 & 29 & 7 & 7 & 18 & 13 & 10 & 6 & 3 & 6 & 3  \\
Problem 30 & 52 & 10000 & 1 & 1 & 14 & 9 & 23 & 14 & 11 & 5 & 1 & 1  \\
Problem 31 & 1 & 1 & 31 & 1 & 1 & 1 & 10000 & 1 & 1 & 1 & 1 & 1  \\
Problem 32 & 45 & 9 & 28 & 6 & 6 & 4 & 4 & 7 & 6 & 3 & 5 & 3  \\
Problem 33 & 48 & 119 & 1 & 1 & 10 & 5 & 18 & 11 & 10 & 4 & 1 & 1  \\
Problem 34 & 47 & 12 & 29 & 6 & 8 & 10 & 12 & 8 & 7 & 3 & 5 & 2  \\
Problem 35 & 1 & 10000 & 33 & 10 & 1 & 1 & 3 & 7 & 1 & 1 & 9 & 6  \\
Problem 36 & 5 & 10000 & 4 & 4 & 5 & 7 & 17 & 33 & 6 & 6 & 4 & 4  \\
Problem 37 & 47 & 6 & 30 & 5 & 7 & 19 & 6 & 9 & 5 & 3 & 4 & 3  \\
Problem 38 & 2 & 1 & 32 & 1 & 1 & 1 & 10000 & 1 & 1 & 1 & 1 & 1  \\
Problem 39 & 49 & 1 & 31 & 1 & 1 & 1 & 10000 & 1 & 1 & 1 & 1 & 1  \\
Problem 40 & 51 & 123 & 32 & 9 & 10 & 7 & 23 & 11 & 12 & 5 & 9 & 4  \\
Problem 41 & 43 & 3790 & 28 & 9 & 12 & 7 & 14 & 18 & 9 & 5 & 8 & 5  \\
Problem 42 & 1 & 1 & 10000 & 1 & 1 & 1 & 10000 & 1 & 1 & 1 & 1 & 1  \\
Problem 43 & 49 & 26 & 31 & 8 & 10 & 5 & 14 & 9 & 8 & 4 & 7 & 4  \\
Problem 44 & 47 & 11 & 28 & 6 & 8 & 7 & 10 & 8 & 7 & 2 & 5 & 3  \\
Problem 45 & 46 & 8 & 31 & 6 & 6 & 6 & 9 & 7 & 5 & 2 & 5 & 2  \\
Problem 46 & 47 & 10 & 28 & 5 & 7 & 4 & 4 & 7 & 6 & 3 & 6 & 3  \\
Problem 47 & 48 & 8 & 26 & 9 & 9 & 31 & 16 & 12 & 8 & 4 & 7 & 3  \\
Problem 48 & 44 & 187 & 27 & 8 & 9 & 6 & 20 & 13 & 9 & 4 & 6 & 4  \\
\hline
\end{tabular}}
\end{table}
\FloatBarrier

\subsection{Results Analysis}

The results of the comprehensive 48-problem evaluation yield several important findings:

\begin{enumerate}
    \item \textbf{Hybrid Algorithm Superiority:} All proposed hybrid algorithms outperform the traditional baselines. The best performer, \textbf{Optimized BFMS}, achieved an average CPU time of \textbf{10.25 $\mu s$} and secured wins on \textbf{32 out of 48 problems} (66.7\%). It achieved a substantial \textbf{6.01x speedup} over Bisection ($61.64 \mu s$) and a massive \textbf{73.9x speedup} over the slowest baseline, False Position ($757.44 \mu s$). The \textbf{Optimized TFMS} algorithm ranked second with an average CPU time of \textbf{12.45 $\mu s$}, followed closely by \textbf{Optimized BF} at \textbf{12.65 $\mu s$}.
    
    \item \textbf{Comparison with Modern Methods:} The proposed hybrids consistently outperformed state-of-the-art techniques. The best hybrid, \textbf{Opt.BFMS}, proved to be \textbf{1.75 times faster} than Brent's method (SciPy, $17.89 \mu s$) and \textbf{2.46 times faster} than Quadratic Interval Refinement (QIR, $25.25 \mu s$). This highlights the significant performance advantage of the proposed multi-phase architecture.
\end{enumerate}

\subsection{Statistical and Significance Analysis}

The experimental design employed a high-volume single-loop strategy ($N=10,000$), yielding robust statistics. The coefficient of variation (CV) for the CPU time measurements is consistent with high-performance micro-benchmarking (e.g., for Opt.BFMS: $\text{CV} = \frac{10.86}{10.25} \approx 1.06$).

The performance gains are statistically significant. Comparing the CPU times of the best hybrid, \textbf{Optimized BFMS} (mean = $10.25 \mu s$), against the most reliable traditional algorithm, \textbf{Bisection} (mean = $61.64 \mu s$), confirms a clear separation in performance distributions.

\subsection{Comparison with Traditional Methods}

The hybrid algorithms demonstrate unmatched superiority over traditional methods across all performance criteria. Table~\ref{tab:traditional_comparison} provides a focused comparison with the fastest traditional baseline.

\begin{table}[H]
\centering
\captionof{table}{Performance Comparison: Hybrid vs Traditional Methods. This table focuses on the performance gap between the proposed best hybrid algorithm and the best-performing traditional numerical method, highlighting the substantial improvements achieved.}
\label{tab:traditional_comparison}
\begin{tabular}{lccc}
\hline
Metric & Best Hybrid & Best Traditional & Improvement \\
\hline
Avg CPU Time (us) & 10.25237277 & 61.64110822 & \textbf{6.01x faster} \\
Success Rate & 100\% & 100\% & Equal reliability \\
\hline
\end{tabular}
\end{table}

\textbf{Operational Performance:} With an average CPU time \textbf{6.0 times faster} than Bisection, the hybrid approach offers a significant advantage for real-time applications where computational resources are constrained or high-throughput root finding is required.

\textbf{Reliability:} Despite the increased speed, the hybrid algorithms maintain a \textbf{100\% success rate} across the test suite, preserving the robust convergence guarantee of traditional bracketing methods.

\subsection{In Comparison with State-of-the-Art Methods} 

These hybrid algorithms can also surpass the state-of-the-art modern techniques such as the QIR methods and the implementations of the Brent method. See Table~\ref{tab:state_of_art_comparison}.

\begin{table}[H]
\centering
\captionof{table}{Performance Comparison: Hybrid vs State-of-the-Art Methods. This table compares the proposed hybrid algorithms with modern methods including QIR and two variants of Brent's method (SciPy and Implementation), demonstrating the continued superiority of our approach.}
\label{tab:state_of_art_comparison}
\begin{tabular}{lcccc}
\hline
Metric & Best Hybrid & QIR & Brent (SciPy) & Brent (Impl.) \\
\hline
Avg CPU Time ($\mu s$) & 10.25 & 25.25 & 17.89 & 18.10 \\
Wins & 32 & 1 & 1 & 3 \\
Win Percentage & 66.67\% & 2.08\% & 2.08\% & 6.25\% \\
Speedup vs QIR & 2.46x faster & - & 1.41x faster & 1.39x faster \\
Speedup vs Brent (SciPy) & 1.75x faster & 1.41x slower & - & 1.01x slower \\
\hline
\end{tabular}
\end{table}

\subsection{Analysis vs. State-of-the-Art}

A key finding of this study is the consistent outperformance of our top hybrid, Opt.BFMS, against established methods.

\textbf{Superiority over QIR:} The proposed Opt.BFMS algorithm significantly outperforms Quadratic Interval Refinement, achieving an average CPU time of \textbf{10.25 $\mu s$} compared to QIR's \textbf{25.25 $\mu s$} (a 2.46x speedup). Furthermore, the hybrid method demonstrates superior consistency, winning on 32 problems compared to QIR's single win.

\textbf{Competitive Edge over Brent:} While Brent's method (SciPy) remains a competitive standard ($17.89 \mu s$), Opt.BFMS achieves a \textbf{1.75x speedup}. The hybrid algorithm's simpler, deterministic architecture avoids the overhead of Brent's complex decision trees while effectively exploiting similar interpolation opportunities.

\subsection{Algorithm-Specific Performance Analysis}

Each of the hybrid algorithms underpinned by the following distinctions in performance characteristics has contributed to the general description of the effectiveness of the approach:

\textbf{Opt.BFMS (Best Overall):} This algorithm achieves the optimal balance of speed and reliability, with the lowest average CPU time ($10.25 \mu s$). The three-phase structure (Bisection $\to$ False Position $\to$ Modified Secant) effectively extracts maximum information from each function evaluation.

\textbf{Opt.TFMS (Second Best):} Ranking second with an average time of \textbf{12.45 $\mu s$}, this method demonstrates the viability of trisection-based hybrids. While slightly slower than Opt.BFMS on average, it occasionally outperforms others on specific problem types.

\textbf{Opt.BF (Strong Baseline):} With an average time of \textbf{12.65 $\mu s$}, this two-phase method proves that even without the secant step, a simple Bisection-False Position hybrid is highly effective, significantly outperforming traditional methods.

\textbf{Opt.TF:} Averaging \textbf{15.02 $\mu s$}, this method serves as an effective baseline for trisection hybrids. While it lags behind the secant-augmented variants, it remains robust and competitive against non-hybrid methods.

\begin{figure}
    \centering
    \includegraphics[width=0.8\linewidth]{avg_cpu_time_per_algorithm.png}
    \captionof{figure}{Average CPU time (in microseconds) for each algorithm across all test problems. Lower values indicate higher computational efficiency.}
    \label{fig:avg_cpu_time}
\end{figure} 

\begin{figure}
    \centering
    \includegraphics[width=0.8\linewidth]{cpu_time_lineplot_per_problem.png}
    \captionof{figure}{CPU time per test problem for each algorithm. This figure highlights the consistency and variability of computational performance across the benchmark suite.}
    \label{fig:cpu_time_per_problem}
\end{figure}

\begin{figure}
    \centering
    \includegraphics[width=0.8\linewidth]{fastest_algorithm_per_problem_lineplot.png}
    \captionof{figure}{Fastest algorithm for each test problem, as determined by minimum CPU time. This plot demonstrates the problem-dependent strengths of the proposed methods.}
    \label{fig:fastest_algorithm}
\end{figure}

\begin{figure}
    \centering
    \includegraphics[width=0.8\linewidth]{avg_cpu_time_lineplot_overall.png}
    \captionof{figure}{Overall average CPU time trend across all problems and algorithms. This figure provides a holistic view of computational efficiency.}
    \label{fig:avg_cpu_time_overall}
\end{figure}

\begin{figure}
    \centering
    \includegraphics[width=0.8\linewidth]{top3_fastest_algorithms_lineplot.png}
    \captionof{figure}{Top three fastest algorithms per problem, ranked by CPU time. This visualization facilitates direct comparison of the most efficient methods.}
    \label{fig:top3_fastest}
\end{figure}
\FloatBarrier

\section{Discussion}\label{sec:discussion}

It was by experimentation as well as theory that the work proved that these multi-phase jointly hybrid algorithms that work one after the other or sequentially have a winning combination in speed, efficiency, and robustness. The first-ranked algorithm, namely \textbf{Opt.BFMS}, proved to have the least average CPU time along with the very minimum average NFE values, confirming the effectiveness of the algorithm in terms of improvement of performance being grounded on algorithm efficiency. Interpretation of these findings with respect to current methods, as well as a look at other broader impacting issues, traditionally follows.

Three phases architecture maximized the success of Opt.BFMS. The first measure is a bisection that guarantees interval reduction: thus, a solid foundation for the subsequent better speedup steps. Then, a false position and modified secant steps permit powerful superlinear convergence to quickly sharpen the location of the root. Empirical data indicate that the overheads of doing these three different steps each iteration more than compensate for the striking reduction in the total number of iterations required for convergence.

The trisection-based algorithms (Opt.TF and Opt.TFMS) were relatively efficient in terms of NFE. For instance, among all algorithms, Opt.TFMS turned out to be the second-most efficient algorithm. However, their somewhat longer CPU times compared to the bisection-based counterparts iowy suggest that, while the time taken by evaluating two interior points for trisection can outweigh the gain from more aggressive interval reductions at times.

The experimental framework presented reliable and reproducible results thanks to the combination of SymPy for symbolic computation and SQLite for data management. It was found that for every algorithm, 10000 repetitions of the experiment made the results statistically valid, whereas tolerance of very high precision equal to $10^{-14}$ ensured root approximations to be quite accurate. The all-embracing test suite has transcendental, polynomial, trigonometric, logarithmic, and exponential functions included, thereby demonstrating the versatility of the algorithms across mathematics.

\section{Conclusion}\label{sec:conclusion}

This document presents four new hybrid algorithms bracketing roots. In these concepts, the bracketing algorithms are then combined with traditional bracketing methods with exceleration in interpolation techniques. All the hybrid algorithms enter upon more accelerated convergence over traditional methods, out of which Opt.BFMS stands first out of all four. The algorithms were achieved equally in all test functions. It also shows a significant improvement of CPU time over classical and modern algorithms. Clearly, these results point out the efficiency of multi-phase hybrid methods in root finding numerically and provide an encouragement for immediate use in many disciplines-science and engineering-related problems.

Their contributions to theory involve the development of four distinct hybrid algorithms that integrate bracketing techniques and open ones, while ensuring their convergence, just that it is faster. The experimental analysis, with its large diversity of 48 test functions and a very rigid high-precisions framework, attested the statistical significance of the results and their robustness. 

The algorithms are competitive against a wide spectrum of mathematical test problems, but there is still a recommendation to test them even in a real-world engineering and physics application case for establishing generality. Future works will include expanding the testing suite by adding real-life case studies and going beyond in hybridization strategies.

\section{Declarations}\label{sec:declaration}

\subsection{Funding}
The author did not receive support from any organization for the submitted work. No funding was received to assist with the preparation of this manuscript. No funding was received for conducting this study. No funds, grants, or other support was received.

\subsection{Competing interests}
The author has no relevant financial or non-financial interests to disclose. The author has no competing interests to declare that are relevant to the content of this article. The author certifies that they have no affiliations with or involvement in any organization or entity with any financial interest or non-financial interest in the subject matter or materials discussed in this manuscript. The author has no financial or proprietary interests in any material discussed in this article.

\subsection{Ethics approval}
Not applicable. This study is purely computational and does not involve human participants or animals.

\subsection{Data availability}
The datasets generated and/or analyzed during the current study are available from the corresponding author on reasonable request. The SymPy implementations of all algorithms and experimental scripts are available at: \url{https://github.com/Abdelrhman-Ellithy/NM-Algorithms/}

\subsection{Authors’ contributions}
The author conceived the study, designed the algorithms, performed the experiments and analysis, and wrote the manuscript.

\begin{thebibliography}{99}

\bibitem{sabharwal2019blended}
Sabharwal, C.L. Blended Root Finding Algorithm Outperforms Bisection and Regula Falsi Algorithms. \textit{ Mathematics 2019, 7, 1118}.
    
\bibitem{sabharwal2019hybrid}
Sabharwal, C. L., and Aggarwal, S. (2019). Hybrid Algorithm Improving Bisection, Regula Falsi, Dekker, Brent Algorithms. \textit{International Journal of Latest Research in Engineering and Technology (IJLRET)}, \url{http://www.ijlret.com/Papers/Vol-05-issue-12/1.B2019072.pdf}.

\bibitem{badr2022novel}
Badr, E., Attiya, H., and El Ghamry, A. (2022). Novel hybrid algorithms for root determining using advantages of open methods and bracketing methods. \textit{Alexandria Engineering Journal}, 61(12), 11579--11588.
  
\bibitem{sabharwal2021iterative}
Sabharwal, C. L., and Aggarwal, S. (2021). An Iterative Hybrid Algorithm for Roots of Non-Linear Equations. \textit{Engineering Proceedings}, 2(1), 7.

\bibitem{badr2021comparative}
Badr, E., Almotairi, S., and El Ghamry, A. (2021). A Comparative Study among New Hybrid Root Finding Algorithms and Traditional Methods. \textit{Mathematics}, 9(11), 1306.

\bibitem{demir2008trisection}
Demir, A. (2008). Trisection method by k-Lucas numbers. \textit{Applied Mathematics and Computation}, 198(1), 339--345.

\bibitem{kelley1995iterative}
Kelley, C. T. (1995). \textit{Iterative Methods for Linear and Nonlinear Equations}. Society for Industrial and Applied Mathematics (SIAM).
    
\bibitem{brent1973algorithms}
Brent, R. P. (1973). \textit{Algorithms for Minimization Without Derivatives}. Prentice-Hall.
                    
\bibitem{abbott2014quadratic}
Abbott, J. (2014). Quadratic interval refinement for real roots. \textit{ACM Communications in Computer Algebra}, 48(1/2), 3--12.
    
\bibitem{kerber2015root}
Kerber, M., and Sagraloff, M. (2015). Root refinement for real polynomials using quadratic interval refinement. \textit{Journal of Computational and Applied Mathematics}, 280, 377--395.
    
\bibitem{dowell1972modified}
Dowell, M., and Jarratt, P. (1972). A modified regula falsi method for computing the root of an equation. \textit{BIT Numerical Mathematics}, 12(4), 503--508.
    
\bibitem{dekker1969finding}
Dekker, T. J. (1969). Finding a zero by means of successive linear interpolation. In B. Dejon \& P. Henrici (Eds.), \textit{Constructive Aspects of the Fundamental Theorem of Algebra} (pp. 37-48). Wiley-Interscience.

\end{thebibliography}

\end{document} 